% Source: MQ18a_v1.html

\documentclass[11pt]{article}
\usepackage[margin=1in]{geometry}
\usepackage{amsmath}
\usepackage{amssymb}
\usepackage{siunitx}
\usepackage{enumitem}

\title{MQ18a: Heat and Calorimetry}
\author{}
\date{}

\begin{document}
\maketitle

\section*{MQ18a: Heat and Calorimetry}

\begin{enumerate}[leftmargin=*,label=\textbf{Question \arabic*.}]

\item \hfill \textit{1 pts}

Calculate the volume in liters of 1.8 moles of an ideal gas at STP (1 atm and 273 K). (1 L = 10$^{-3}$ m$^{3}$)



\emph{V} = \rule{2.5cm}{0.4pt} L  (Note the units!)

\item \hfill \textit{1 pts}

A 7-L tank at 30°C contains 27.77 × 10$^{-4}$ kg of element 4 listed in the table below. Calculate the tank pressure in kPa. (1 L = 10$^{-3}$ m$^{3}$)

Element
Molar Mass (g/mol)

1
Helium
4.00

2
Neon
20.2

3
Argon
40.0

4
Krypton
83.8



\emph{p} = \rule{2.5cm}{0.4pt} kPa  (Note the units!)

\item \hfill \textit{1 pts}

A small bubble forms at a depth of 8 m in a freshwater lake. If the bubble's initial volume is 158 mm$^{3}$ and it expands without change in temperature as it rises, what is its volume in mm$^{3}$ just as it reaches the surface? Assume the pressure at the surface is 94,043 Pa and the water's density is 1000 kg/m$^{3}$.



\emph{V} = \rule{2.5cm}{0.4pt} mm$^{3}$

\item \hfill \textit{1 pts}

Calculate the number of moles in 81 g of the molecule C$_{2}$H$_{4}$.

\item \hfill \textit{1 pts}

Calculate the mass of one molecule of C$_{2}$H$_{3}$OH in kg and supply the missing prefix below.



\emph{m}= \rule{2.5cm}{0.4pt} × 10$^{-27}$ kg

\end{enumerate}

\end{document}
